% !PRINTED 430
\weberpage{430}
We must confine ourselves here to a brief statement of the results, and shall make no applications of them in what follows, so that this exposition, which can be fully intelligible only to a reader somewhat familiar with the theory of transcendental functions, may also be passed over without disturbing the continuity.

Let $\omega$ denote a variable whose imaginary part is positive, so that
\weq{1}{q = e^{\pi i \omega}}
is a quantity whose absolute value is a proper fraction. We shall consider the following three functions of this variable $\omega$:
\weq{2}{\eta(\omega) = q^{1/12}\prod_{1,\infty}^{\nu}(1 - q^{2\nu}) = q^{1/12}\sum_{-\infty,\infty}^{\nu}(-1)^\nu q^{3\nu^2+\nu} = \sum_{-\infty,\infty}^{\nu}(-1)^\nu q^{\frac{(6\nu+1)^2}{12}},}
\weq{3}{f(\omega) = q^{-\frac{1}{24}}\prod(1 + q^{2\nu-1}),}
\weq{4}{\gamma_2(\omega) = \frac{f(\omega)^{24} - 16}{f(\omega)^8}\footnote{Cf.: \gesperrt{Weber}, ``Elliptische Functionen und algebraische Zahlen''. Braunschweig 1891. S. 63, 86, 149.}.}
In the transformation theory of the elliptic functions\footnote{Ibid., S. 263, 316.} an equation of degree 6 is derived, whose roots are the six quantities
\weq{5}{v_\infty = 5\left(\frac{\eta(5\omega)}{\eta(\omega)}\right)^2, \qquad v_\xi = \left(\frac{\eta\left(\dfrac{-24\xi+\omega}{5}\right)}{\eta(\omega)}\right)^2}
when $\xi$ runs through a complete system of residues to the modulus 5. This equation reads
\weq{6}{v^6 + 10v^3 - \gamma_2(\omega)v + 5 = 0,}
and, as one sees, has great similarity to the resolvent of degree 6 considered in the preceding section
\weq{7}{U^6 - 10zU^3 + z^2U + 5z^2 = 0.}
These two equations pass directly into one another, if one puts
\weq{8}{z = \gamma_2(\omega)^3}
\weq{9}{U = -\gamma_2(\omega)v}
%CONT>>
